\chapter{Multiplicando números negativos}\label{ch:g8:negprod}

O Ano 7 somou e subtraiu \omterm{def:g7:negatives:def}{números relativos} (\cref{ch:g7:negatives});
falta multiplicá-los e dividi-los. Uma pergunta domina o capítulo: por
que ``menos vezes menos'' deveria ser mais? Damos uma razão, e não
apenas uma regra.

\section{As regras de sinal}

\begin{theorem}[Sinal de um produto]\label{thm:g8:negprod:rules}
O \omterm{def:g2:mult:def}{produto} de dois \omterm{def:g7:negatives:def}{números relativos} tem distância até o \omterm{def:g1:counting:zero}{zero} igual ao
\omterm{def:g2:mult:def}{produto} das distâncias, e o sinal dele é dado por:
\begin{center}
\renewcommand{\arraystretch}{1.3}
\begin{tabular}{c|cc}
$\times$ & $+$ & $-$ \\
\hline
$+$ & $+$ & $-$ \\
$-$ & $-$ & $+$ \\
\end{tabular}
\end{center}
\emph{Sinais iguais: \omterm{def:g2:mult:def}{produto} positivo. Sinais opostos: \omterm{def:g2:mult:def}{produto}
negativo.} A mesma tabela vale para os \omterm{def:g3:division:remainder}{quocientes}.
\end{theorem}

\begin{proof}[Por que $(-1) \times (-1) = +1$]
Primeiro, $3 \times (-4)$ quer dizer $(-4) + (-4) + (-4) = -12$: um
positivo vezes um negativo é negativo. Agora observe o padrão:
\[
3 \times (-4) = -12, \quad
2 \times (-4) = -8, \quad
1 \times (-4) = -4, \quad
0 \times (-4) = 0 .
\]
A cada passo, o primeiro fator cai $1$ e o resultado \emph{sobe} $4$.
Continuando o padrão mais um passo:
\[
(-1) \times (-4) = 0 + 4 = +4 .
\]
Qualquer outra escolha quebraria a regularidade da aritmética (a
distributividade). Logo menos vezes menos tem de ser mais.
\end{proof}

\begin{omfigure}
\begin{tikzpicture}
\begin{axis}[omaxis,
  width=7.8cm, height=5.4cm,
  xmin=-2.6, xmax=3.6, ymin=-14, ymax=10,
  xlabel={$k$}, ylabel={$k \times (-4)$},
  xtick={-2,-1,1,2,3}, ytick={-12,-8,-4,4,8}]
  \addplot[omDef, thick, domain=-2.3:3.3] {-4*x};
  \addplot[omDef, only marks, mark=*, mark size=1.6pt] coordinates
    {(3,-12) (2,-8) (1,-4) (0,0) (-1,4) (-2,8)};
\end{axis}
\end{tikzpicture}

{\small O argumento do padrão numa imagem: os pontos
$(k,\ k \times (-4))$ se alinham, e continuar a reta além de $k = 0$
obriga $(-1)\times(-4) = +4$ e $(-2)\times(-4) = +8$.}
\end{omfigure}

\begin{example}\label{ex:g8:negprod:products}
\begin{align*}
(-7) \times (+6) &= -42 && \text{(sinais opostos)} \\
(-5) \times (-8) &= +40 && \text{(sinais iguais)} \\
(+2.5) \times (-4) &= -10 \\
(-36) \div (-9) &= +4 \\
(+15) \div (-2) &= -7.5 .
\end{align*}
\end{example}

\section{Produtos de vários fatores}

\begin{proposition}[Sinal de um produto longo]\label{prop:g8:negprod:many}
O sinal de um \omterm{def:g2:mult:def}{produto} de vários fatores não nulos depende apenas da
quantidade de fatores negativos:
\begin{itemize}
  \item uma quantidade \emph{par} de fatores negativos dá \omterm{def:g2:mult:def}{produto}
        positivo;
  \item uma quantidade \emph{\omterm{def:g3:numbers:evenodd}{ímpar}} dá \omterm{def:g2:mult:def}{produto} negativo.
\end{itemize}
\end{proposition}

\begin{proof}
Os fatores negativos se emparelham, e cada par contribui com sinal
positivo (\cref{thm:g8:negprod:rules}); uma quantidade \omterm{def:g3:numbers:evenodd}{ímpar} deixa um
fator negativo sem par, o que torna o \omterm{def:g2:mult:def}{produto} inteiro negativo.
\end{proof}

\begin{example}\label{ex:g8:negprod:many}
$(-2) \times (+3) \times (-5) \times (-1)$: três fatores negativos
(\omterm{def:g3:numbers:evenodd}{ímpar}), então o \omterm{def:g2:mult:def}{produto} é negativo; distâncias:
$2 \times 3 \times 5 \times 1 = 30$. Resultado: $-30$. Decida o sinal
\emph{primeiro} e depois multiplique as distâncias --- duas tarefas
fáceis em vez de uma cheia de armadilhas.
\end{example}

\begin{example}[Potências de negativos]\label{ex:g8:negprod:powers}
$(-2)^2 = (-2) \times (-2) = +4$, e
$(-2)^3 = (-2)^2 \times (-2) = -8$: expoentes pares dão valores
positivos, e expoentes \omterm{def:g3:numbers:evenodd}{ímpares} mantêm o sinal. Cuidado com a notação:
$(-3)^2 = 9$, mas $-3^2 = -(3 \times 3) = -9$.
\end{example}

\section{Calcular com as quatro operações}

\begin{method}[Cálculo seguro com números relativos]\label{met:g8:negprod:compute}
\begin{enumerate}
  \item Respeite as prioridades do \cref{ch:g7:priorities}: parênteses,
        depois $\times$ e $\div$, depois $+$ e $-$;
  \item em cada multiplicação ou divisão, descubra \emph{primeiro o
        sinal} e depois a distância;
  \item reescreva uma etapa por linha.
\end{enumerate}
\end{method}

\begin{example}\label{ex:g8:negprod:compute}
\begin{align*}
5 - 3 \times (-4)
&= 5 - (-12) && \text{(o produto primeiro: sinais opostos)}\\
&= 5 + 12 = 17 .
\end{align*}
\begin{align*}
\frac{(-6) \times (-10)}{-4}
&= \frac{+60}{-4} && \text{(numerador: sinais iguais)}\\
&= -15 && \text{(quociente: sinais opostos).}
\end{align*}
\end{example}

\begin{example}[Substituir valores negativos]\label{ex:g8:negprod:substitute}
Calcule $E = 3x^2 - 2x$ para $x = -5$, com parênteses em torno do
valor:
\[
E = 3 \times (-5)^2 - 2 \times (-5)
= 3 \times 25 - (-10)
= 75 + 10 = 85 .
\]
Os parênteses em $(-5)^2$ são essenciais: sem eles, o quadrado se
aplicaria só ao $5$.
\end{example}

\section{Exercícios}

\begin{exercise}[$\star$]\label{exo:g8:negprod:1}
Calcule:
\[
(-8) \times (+7), \qquad
(-9) \times (-6), \qquad
(+12) \times (-0.5), \qquad
(-1) \times (-1) \times (-1).
\]
\end{exercise}

\begin{exercise}[$\star$]\label{exo:g8:negprod:2}
Calcule:
\[
(-63) \div (+9), \qquad
(-8) \div (-16), \qquad
\frac{45}{-5}, \qquad
\frac{-7.2}{-8} .
\]
\end{exercise}

\begin{exercise}[$\star$]\label{exo:g8:negprod:3}
Dê apenas o \emph{sinal} de cada \omterm{def:g2:mult:def}{produto}, sem calculá-lo:
\[
(-3) \times (-7) \times (+2) \times (-5); \qquad
(-1)^{10}; \qquad
(-2)^{7}; \qquad
(-4) \times 0 \times (-6).
\]
\end{exercise}

\begin{exercise}[$\star$]\label{exo:g8:negprod:4}
Calcule:
\[
(-2)^4, \qquad
-2^4, \qquad
(-3)^3, \qquad
(-1)^{2026} .
\]
\end{exercise}

\begin{exercise}[$\star$]\label{exo:g8:negprod:5}
Calcule passo a passo, respeitando as prioridades:
\[
7 + 2 \times (-6), \qquad
(-4) \times (5 - 9), \qquad
18 \div (-3) - (-2) .
\]
\end{exercise}

\begin{exercise}[$\star$]\label{exo:g8:negprod:6}
Calcule:
\[
\frac{(-5) \times (+8)}{-4}, \qquad
\frac{(-3) \times (-6)}{(-2) \times (+9)} .
\]
\end{exercise}

\begin{exercise}[$\star$]\label{exo:g8:negprod:7}
Calcule o valor para $x = -3$:
\[
5x, \qquad x^2, \qquad -x^2, \qquad 2x^2 + 4x, \qquad (2x)^2 .
\]
\end{exercise}

\begin{exercise}[$\star\star$]\label{exo:g8:negprod:8}
Complete cada igualdade:
\[
(-6) \times \;?\; = 42, \qquad
\;?\; \div (-4) = -2.5, \qquad
(-5) \times \;?\; = -1 .
\]
\end{exercise}

\begin{exercise}[$\star\star$]\label{exo:g8:negprod:9}
Verdadeiro ou falso? Justifique ou dê um contraexemplo.
\begin{enumerate}
  \item O \omterm{def:g2:mult:def}{produto} de dois \omterm{def:g7:negatives:def}{números relativos} é sempre pelo menos tão
        grande quanto a \omterm{def:g1:addition:def}{soma} deles.
  \item O quadrado de um \omterm{def:g7:negatives:def}{número relativo} nunca é negativo.
  \item Se um \omterm{def:g2:mult:def}{produto} de três fatores é positivo, os três fatores são
        positivos.
\end{enumerate}
\end{exercise}

\begin{exercise}[$\star\star$]\label{exo:g8:negprod:10}
Num quiz, cada resposta errada vale $-3$ pontos e cada resposta certa,
$+5$. Zoe respondeu $20$ questões e fez $36$ pontos. Quantas respostas
foram certas? (Teste valores ou monte o cálculo
$5r - 3(20 - r) = 36$.)
\end{exercise}

\begin{exercise}[$\star\star\star$]\label{exo:g8:negprod:11}
Usando a distributividade (como na demonstração da
\cref{thm:g8:negprod:rules}), desenvolva e justifique cada etapa de
\[
0 = (-1) \times 0 = (-1) \times \bigl(1 + (-1)\bigr)
= (-1) \times 1 + (-1) \times (-1),
\]
e conclua que $(-1) \times (-1)$ tem de ser igual a $+1$.
\end{exercise}

\section{Problema: por que menos vezes menos é mais}

\begin{problem}[{Problema de fim de semana --- as regras de sinal
deduzidas da distributividade, e a soma alternada
$1 - 2 + 3 - 4 + \dots$}]
\label{pb:g8:negprod:1}
A demonstração do \cref{thm:g8:negprod:rules} continuou um padrão e
afirmou que ``qualquer outra escolha quebraria a distributividade''.
Este problema torna a afirmação exata: partindo apenas da
distributividade (\cref{thm:g7:priorities:distributivity}), você vai
\emph{demonstrar} as regras de sinal --- sem desenho e sem padrão, do
jeito que a álgebra faz --- e depois pô-las para trabalhar nas
potências de $-1$ e numa \omterm{def:g1:addition:def}{soma} famosa de sinais alternados. Ao longo de
tudo, $a$ e $b$ são \omterm{def:g7:negatives:def}{números relativos}; damos por certas as regras de
adição do \cref{ch:g7:negatives}, que $1 \times a = a$ e que um \omterm{def:g2:mult:def}{produto}
pode ser calculado em qualquer ordem (então a distributividade vale dos
dois lados de um \omterm{def:g2:mult:def}{produto}).

\medskip
\noindent\textbf{Parte I --- As regras de sinal são teoremas.}
\begin{enumerate}
  \item Demonstre que $0 \times a = 0$ para todo \omterm{def:g7:negatives:def}{número relativo} $a$.
        (Escreva $0 = 0 + 0$, desenvolva $(0 + 0) \times a$ e pergunte
        que número, somado a $0 \times a$, devolve $0 \times a$.)
  \item Desenvolva $\bigl(1 + (-1)\bigr) \times a$ para mostrar que
        $(-1) \times a + a = 0$ e conclua que
        \[
        (-1) \times a = -a :
        \]
        multiplicar por $-1$ dá o oposto. Compare com o
        \cref{exo:g8:negprod:11}, que é o caso $a = -1$.
  \item Deduza que $(-a) \times b = -(a \times b)$: um sinal de menos
        sai de um \omterm{def:g2:mult:def}{produto} sem alteração.
  \item Deduza que $(-a) \times (-b) = a \times b$: dois sinais de
        menos se cancelam.
  \item Todo \omterm{def:g7:negatives:def}{número relativo} é a distância dele até o \omterm{def:g1:counting:zero}{zero} com um sinal
        na frente: $a = +d$ ou $a = -d$. Explique como as questões 3 e
        4 demonstram os quatro casos da tabela de sinais da
        \cref{thm:g8:negprod:rules}, distâncias incluídas.
\end{enumerate}

\medskip
\noindent\textbf{Parte II --- Potências de $-1$.}
\begin{enumerate}[resume]
  \item Calcule $(-1)^2$, $(-1)^3$, $(-1)^4$, $(-1)^5$ e depois dê, com
        justificativa pela \cref{prop:g8:negprod:many}, o valor de
        $(-1)^n$ para todo número inteiro $n \geq 1$.
  \item Dê (sem calcular nenhuma distância) o sinal de
        \[
        (-1) \times (-2) \times (-3) \times \dots \times (-10),
        \]
        e escreva a distância desse \omterm{def:g2:mult:def}{produto} como um \omterm{def:g2:mult:def}{produto} de números
        inteiros, sem efetuá-lo.
  \item Mostre que $(-a)^2 = a^2$ e que $(-a)^3 = -a^3$. Para que
        expoentes o sinal de menos sobrevive?
  \item Use as regras de sinal para mostrar que o quadrado de um \omterm{def:g7:negatives:def}{número
        relativo} nunca é negativo e deduza que nenhum \omterm{def:g7:negatives:def}{número relativo}
        $x$ satisfaz $x^2 = -4$.
  \item Zoe afirma: ``$(-1)^{2026} + (-1)^{2027} = 0$ e, de modo mais
        geral, duas potências consecutivas de $-1$ sempre se
        cancelam''. Ela tem razão? Justifique.
\end{enumerate}

\medskip
\noindent\textbf{Parte III --- A \omterm{def:g1:addition:def}{soma} alternada.}
Com as regras de sinal garantidas, podemos calcular \omterm{def:g1:addition:def}{somas} que misturam
os dois sinais. Para um número inteiro $n \geq 1$, seja
\[
A_n = 1 - 2 + 3 - 4 + \dots
\]
a \omterm{def:g1:addition:def}{soma} dos números inteiros de $1$ a $n$ com sinais alternados: o
último termo é $+n$ quando $n$ é \omterm{def:g3:numbers:evenodd}{ímpar} e $-n$ quando $n$ é par. Assim,
$A_1 = 1$ e $A_2 = 1 - 2 = -1$.
\begin{enumerate}[resume]
  \item Calcule $A_3$, $A_4$, $A_5$, $A_6$ e $A_7$. O que você
        conjectura?
  \item Suponha $n$ par. Agrupe os termos dois a dois,
        $(1 - 2) + (3 - 4) + \dots$, conte os pares e demonstre que
        $A_n = -\frac{n}{2}$.
  \item Suponha $n$ \omterm{def:g3:numbers:evenodd}{ímpar}. Explique por que $A_n = A_{n-1} + n$ e
        deduza da questão anterior que $A_n = \frac{n + 1}{2}$.
  \item Calcule $A_{100}$, $A_{101}$ e $A_{2026}$.
  \item Alma calculou $A_{101} - A_{100} = 51 - (-50) = 101$ e se
        surpreendeu ao encontrar exatamente $101$. Ela deveria ter se
        surpreendido? Explique em uma frase e dê o valor de
        $A_{2027} - A_{2026}$ sem nenhuma conta a mais.
\end{enumerate}
\end{problem}
